\documentclass[book.tex]{subfiles}

\begin{document}

\chapter{Quantum Harmonic Oscillators}

\label{chap:quantum_harmonic}
\noindent\rule{11cm}{0.4pt} \\
\textbf{Prerequisites:}
\autoref{chap:harmonic},
\autoref{chap:quantum_mechanics}, \autoref{chap:position_and_momentum} \\
\textbf{Difficulty Level:} ** \\
\noindent\rule{11cm}{0.4pt}
\vspace{5mm}

The quantum harmonic oscillator is a quantum mechanical version of the classical oscillator we saw in \autoref{chap:harmonic}. Quantum harmonic oscillators also arise as a foundational building block for systems in condensed matter and quantum field theory.

\section{The Basic Oscillator}
The energy for a classical harmonic spring system was given by
\begin{equation}
    E = \frac{p^2}{2m} + \frac{1}{2}kx^2
\end{equation}
Quantizing this energy, we arrive at the Hamiltonian for a quantum mechanical oscillator.
\begin{equation}
    \hat{H} = \frac{\hat{p}^2}{2m} + \frac{1}{2}k\hat{x}^2
\end{equation}
How should we think about this equation? Proceeding with intuition based on the classical energy, this operator states that as the energy and momentum grows larger, the system has higher and higher energy.

We start by defining some fundamental types for this system. The state space $\mathcal{S}$ should be a suitable Hilbert space $\mathcal{H}$ of wave functions. For this example, we will take $\mathcal{H}$ to be the space of square-integrable functions over $\mathbb{R}^3$, $L^2(\mathbb{R}^3)$

The time evolution of the system is given by Schrodinger's equation. In our case, the equation becomes
\begin{align}
        i \hbar \frac{d}{dt} | \psi(t) \rangle = \left (\frac{\hat{p}^2}{2m} + \frac{1}{2} k \hat{x}^2 \right ) |\psi(t) \rangle
\end{align}

The parameters are 
\begin{align}
\mathcal{W} = \{m, k, |\psi(0)\rangle \}
\end{align}
where $k$ is the harmonic constant and $|\psi(0)\rangle$ is the initial condition of the wave function.

Recall that the velocity operator is defined by
\[ \hat{v} = -\frac{i\hbar}{m} \frac{\partial}{\partial x} \]
We constrain
\[ \|\hat{v}\|_2 \ll c \]
where $\| \hat{v}\|_2$ is the spectral norm of $\hat{v}$.
Hyperparameters $\mathcal{HP}$ are $c$, the speed of light, and noise parameter $\epsilon$.

In the classical quantum mechanical sense, any self-adjoint operator on $\mathcal{H}$ can be an observable. We choose the position operation $\hat{x}$ to be the observable we can measure. Recall that taking a measurement of $\hat{x}$ projects onto the eigenstates of $\hat{x}$ (which are just the points in $\mathbb{R}^3$). We will assume that we have a noise process $\mathrm{Normal}(0, \epsilon)$ on this measurement

\label{phs:theory:moving:particle}

\section{The Vibrational Modes of a Diatomic Molecule}

The vibrational modes of a diatomic molecule can be determined by considering a single particle in a harmonic potential. Consider a diatomic molecule with atomic masses $m_1$
and $m_2$. The covalent bond between the two atoms can be modeled as a harmonic spring with spring constant $k$. If we define $x$ to be the distance of separation between the two atoms, we have the governing equation for the wave-packet as %\textcolor{red}{TODO: Define wave packet here.}

\begin{align}
\hat{H} \psi &= E \psi \\
-\frac{h^2}{8 \pi^2 \mu} \frac{d^2\psi}{dx^2} + \frac{1}{2}kx^2\psi &= E \psi
\end{align}

Where 
\begin{align}
\mu &= \frac{m_1m_2}{m_1+m_2}
\end{align}
Upon solving this equation, we have an infinite number of solutions with discreet energy levels. In general, the $nth$ wave-packet can be described by:

\begin{align}\label{exmp4b}
\psi_n = \left(\frac{1}{n!2^na\sqrt{\pi}}\right)^{1/2} H_n\left(\frac{x}{a}\right) \exp\left(-\frac{x^2}{2a^2}\right)
\end{align}

Where $a^4=\frac{h^2}{4\pi^2\mu k}$ and the function $H_n(u)$ represents the Hermite polynomials as
\begin{align}
H_0(u) &=1\\
H_1(u) &=2u\\
H_2(u) &= 4u^2 - 2 \\
H_3(u) &= 8u^3-12u
\end{align}
for the first four states. Following this procedure, the energy of the $nth$ state can be described as:

\begin{marginfigure}%
\includegraphics[width=\linewidth]{figures/Physics/quantum_mechanics/harmonic_oscillator/harmonicoscillator.png}
\caption{Wave-packet representations for the eigenstates, $n = 0$ to $7$ for the harmonic oscillator. The horizontal axis shows the position $x$. }%\textcolor{red}{TODO: Can we use this?}}
\label{fig:harmonicoscillator}
\end{marginfigure} 

\begin{align}\label{exmp4c}
E_n = \left( n+\frac{1}{2}\right)\frac{h\nu}{2\pi}
\end{align}

Where $\nu=\sqrt{\frac{k}{\mu}}$ and $n=0,1,2,3,4.....$ are the quantum mechanical modes of motion. As in the case of classical mechanics, the characteristic frequency $\nu$ plays an important role in determining the solutions using the quantum mechanical solutions, as shown in Fig.~\ref{fig:harmonicoscillator}.

\section{Particle in a box}

We can use the quantum harmonic oscillator to help us solve more complex models. Consider now a particle with 1-D motion along the $x$ axis in a box of length L from $x=0$ to $x=L$. The external potential is assigned as 
\begin{align}
V(x,t) = \begin{cases} 
0 & \textrm{ for $0 < x < L$} \\
\rightarrow \infty & \textrm{ for $x \geq L$ and $x \leq 0$} 
\end{cases}
\end{align} 
We illustrate $V(x)$ in Fig.~\ref{fig:partinbox}. 

\begin{marginfigure}
    % \centering    \includegraphics{figures/Physics/quantum_mechanics/schrodingers_equation/particle_in_box.png}
    \centering    \includegraphics{figures/Physics/quantum_mechanics/schrodingers_equation/particle_in_box0.png}
    \caption{The potential energy function for a particle in a box. }
    \label{fig:partinbox}
\end{marginfigure}
% \textcolor{red}{TODO: Replace with our own figure.}

The governing equation for the particle inside the box is:

\begin{align}
\hat{H} \psi &= E \psi \\
-\frac{h^2}{8 \pi^2 m} \frac{d^2\psi}{dx^2} &= E \psi
\end{align}

This equation has the boundary conditions $\psi=0$ at $x=0$ and at $x=L$. The equation can be written in the same form as that of a harmonic oscillator as:

\begin{align}\label{exmp3b}
\frac{d^2\psi}{dx^2} + G^2\psi=0
\end{align}

Where $G^2=\frac{8mE \pi^2}{h^2}$. The solution  for this system is given as:

\begin{align}\label{exmp3c}
\psi = C_1\sin(Gx)+ C_2\cos(Gx)
\end{align}

To satisfy the boundary conditions, we must have $C_2=0$. The remaining solution has infinite possibilities because $\sin(n\pi)=0$ for $n=1,2,3,4.......$ The condition for the solution thus results in:

\begin{align}\label{exmp3d}
G_n = \pi\sqrt{\frac{8m E_n}{h^2}} = \frac{n\pi}{L}
\end{align}

This implies:

\begin{align}\label{exmp3e}
E_n = \frac{h^2}{8mL^2}n^2
\end{align}

Using the condition that $|\psi|^2$ has to be normalized, we have 
\begin{align}
C_1 &= \sqrt{\frac{2}{L}}
\end{align}
Hence, the solution for the wave-packet for the $n^{th}$ quantum state of the particle is:

\begin{marginfigure}[8cm]
% \includegraphics[width=\linewidth]{figures/Physics/quantum_mechanics/harmonic_oscillator/schrodingereqnsoln.png}
\includegraphics[width=\linewidth]{figures/Physics/quantum_mechanics/harmonic_oscillator/schrodingereqnsoln0.png}
\caption{Solution for the wave-packets for the first four states, $n=1,2,3,4$, in a one-dimensional particle in a box}
\label{fig:schrodingereqnsoln}
\end{marginfigure}

\begin{equation}\label{exmp3f}
\psi_n(x) = \sqrt{\frac{2}{L}} \sin \left (\frac{n \pi x}{L}\right )
\end{equation}

One can easily extend this to 3-dimensions, which instead of $n$ would result in $n_x,n_y,n_z$. However, what is more important here is to understand the quantization of the energy levels in terms of n. For varying $n$, we get different solutions of the Schr{\"o}dinger equation in 1-dimension, as shown in Fig.~\ref{fig:schrodingereqnsoln}. 


\end{document}
